%!TEX program = lualatex
\documentclass[ngerman,headheight=70pt]{scrartcl}
\usepackage[ngerman]{babel}
\usepackage{blindtext}
\usepackage[utf8]{luainputenc}
\usepackage{scrlayer-scrpage}
\usepackage[T1]{fontenc}
\usepackage{fontspec}
\usepackage{multicol}
\usepackage{enumitem}
\usepackage[german=quotes]{csquotes}
\usepackage[bottom=2cm,footskip=8mm]{geometry}
\setmainfont{Source Sans Pro}
\MakeOuterQuote{"}

\setlength{\parindent}{0pt}
\setlength{\parskip}{6pt}

\pagestyle{plain.scrheadings}
\thispagestyle{scrheadings}
\lohead{\vspace*{-1.5cm}\includegraphics[scale=0.4]{../images/stupa_logo_rot.pdf}}
\rohead{\\Präsidium des\\Studierendenparlaments}

\lofoot[Universität Hamburg · Präsidium des Studierendenparlaments\\\hrule Von-Melle-Park 5 · D-20146 Hamburg]{Universität Hamburg · Präsidium des Studierendenparlaments\\\hrule Von-Melle-Park 5 · D-20146 Hamburg}
\cofoot[]{}
\cohead[\thepage]{}

\renewcommand{\thesubsection}{TOP \arabic{subsection}}
\renewcommand{\thesubsubsection}{\arabic{subsubsection}.}
\setcounter{subsection}{-1}

\begin{document}
    UHH · StuPa-Präsidium · Von-Melle-Park 5 · D-20146 Hamburg

    %%% PAGE ONE (EDIT LATER) %%%%
    % Hier den Titel des Protokolls eintragen
    \section*{Protokoll der 5. Sitzung des Studierendenparlaments vom 16. Juni 2016}

    % Protokollanten eintragen
    \textbf{Protokoll: Gunhild Berdal}\\
    \textbf{Ort: S27 in der ehemaligen HWP}\\
    \textbf{Beginn: 18.30 Uhr}\\
    \textbf{Ende: 0.50 Uhr}

    % Wohl die komplizierteste Stelle im Protokoll
    % Als Faustregel sollten nicht mehr als 3 Namen auf einer Zeile stehen
    % Einfach an dieser Vorlage orientieren
    % Im Idealfall ist hier nicht viel zu ändern, da kaum jemand fehlt.
    % Rücktritte müssen natürlich pro Protokoll angepasst werden
    \vspace{0.5cm}
    \begin{tabular}{ll}
        Anwesend: & \\
            RCDS (5 Sitze): & Ramon Weilinger \\
             CampusGrün (14 Sitze): & Laura Franzen, Geoffrey Youett, Elena Rysikova, \\
                                    & Philipp Droll,  Melf Johannsen, Tahnee Herzig \\
                                    & Mario Moldenhauer, Christian Lagod, Jim Martens, \\
                                    & Svenja Horn, Mirzo Ulugbek Khatamov, Martin Sievert, \\
                                    & Armin Günther \\
             BGZM-Liste (2 Sitze) : &  \\
             WiWi (2 Sitze): & \\
             Unicorns (5 Sitze): & Johannes Peplow, Annkathrin Löffler, Andreas Hartkamp, \\
                                 & Marielle Hermstrüwer \\
             Liste LINKS (3 Sitze): & Gunhild Berdal, Till Petersen, Sinah Mielich \\
             HWP (2 Sitze): &  \\
             MIN (4 Sitze): &   \\
             SDS* (3 Sitze): & Mena Winkler, Jacob Petersein, Artur Brückmann \\
             Bart-LISTE (2 Sitze): & Timo Zeimet, Dominic Laumer \\
             LHG (1 Sitz): &  \\
             harte zeiten (1 Sitz): & Tobias Berking \\
             Jura (1 Sitz): & \\
             AL (2 Sitze): & Karima Schulze, Henri Weber \\
            & \\
        Entschuldigt: & Antonia Niecke, Ramin Shakiba, Jennifer Maack (RCDS), \\
                      & Yasemin Günther (CampusGrün), Jakob Pape, Karen Martirosian (BGZM), \\
                      & Katharina Kucza (Unicorns), Ailina Salten, Lotte Rullkötter, \\
                      & Jan Detampel (MIN) \\
                      &\\
        Unentschuldigt abwesend: & Benjamin Welling (RCDS), Claas-Friso Hente, Kay Zöllmer (WiWi), \\
                                & Samet Gunay, Ajdina Karahasan (HWP), Thea Wahlers (MIN), \\
                                & Tobias Heisig (LHG), Johann Baumhoefener (Jura)\\
                                &\\
        Rücktritte: & \\
    \end{tabular}
    \newpage
    %%% Vorgeschlagene Tagesordnung %%%
    \underline{Vorgeschlagene Tagesordnung}
    \begin{enumerate}[label={\textbf{Top \theenumi}},leftmargin=*]
        % folgende Zeile bei Sitzungsfortsetzungen einkommentieren und x durch
        % die Zahl des zuletzt behandelten TOPs ersetzen
        %\setcounter{enumi}{x}
        \item
            \begin{enumerate}
                \item Haushalt 2016/2017 1. Lesung (V1617-025) (90 Min.)
                \item Nachtragshaushalt 2015/2016 1. Lesung (V1617-026) (30 Min.)
            \end{enumerate}
        \item Antrag auf Beitritt zum fzs (V1617-027)
        \item Akkreditierung (V1617-024) (60 Min.)
        \item Prüfauftrag Wahl-O-Mat bei StuPa-Wahlen (V1617-022) (30 Min.)
        \item Verlegung der StuPa-Wahl in das Sommersemester (V1617-015) (60 Min.)
        \item Verschiedenes
    \end{enumerate}

    \newpage

    % Hier beginnt Teil A. Bei Fortsetzungen kommt hier nur der Teil über
    % die Feststellung der endgültigen Fassung der Tagesordnung.
    % Für sonstige Sitzungen ist dies bereits das Gerüst für die Formalia.
    % Man muss nur noch die protokollierten Inhalte eintragen.

    {\Large \underline{Teil A}}

    \subsection{Formalia}

    \subsubsection{Geschäftsbericht Präsidium}

    Ramon öffnet die Sitzung um 18.30 Uhr und entschuldigt sich zugleich für die
    Zeit ab 20.30 Uhr. Aus "aktuellem Anlass" (Deutschlandspiel EM) werde Melf
    ihn ab 20.30 Uhr vertreten.

    Er berichtet, dass das Präsidium für die Diskussion mit Frau Fegebank drei
    Themen verabredet hat: Hochschulfinanzierung, Exzellenzinitiative und zivile
    Forschung.

    \subsubsection{Anfragen an das Präsidium}

    Frage: Hat das Präsidium diskutiert, ob alle Präsidiums-Mitglieder auf der
    StuPa-Sitzung anwesend sein müssen?

    Antwort: Die Frage habe das Präsidium nicht ernsthaft erörtert.

    Anmerkung aus dem Parlament: Zukünftig mögen alle Präsidiumsmitglieder in
    Zivil erscheinen.

    \subsubsection{Geschäftsbericht AStA}

    \textbf{Vorstand} (Franziska, Timo):
    \begin{itemize}
        \item Die Demo "Uni für alle statt Exzellenz-Initiative" war mit knapp
        200 Teilnehmer*innen klein, aber fein. Es gab viele Gespräche und
        Unterschriften. Es wird ein Folgetreffen der Aktiven am Dienstag, den
        21.6. geben.
        \item In Auswertung der letzten StuPa-Sitzung und des "Tags der Bundeswehr"
        will der AStA die Diskussion mit den Studierenden der Bundeswehruni
        fortsetzen.
        \item Es gibt am 18. Juni eine Demo gegen Rassismus und Polizeigewalt.
        Ein Kleindealer namens Jaja Diabi ist im Haft gestorben, die Behörden
        behaupten Selbstmord, die Angehörigen bezweifeln es. Auch Prof. Dr.
        Püschel vom UKE ist involviert. Der AStA unterstützt die Demo und will
        mit den progressiven Kräften im UKE in Kontakt bleiben.
        \item Im Akademischen Senat sei die Begehung des Wissenschaftsrats an
        der Uni kritisch ausgewertet worden.
        \item Der Haushalt der VS ist fertiggestellt.
    \end{itemize}

    \textbf{Antidiskriminierung} (Karima):

    Das Referat arbeitet zur Sichtbarmachung sexueller Gewalt (z.B. sexuelle
    Belästigung auf den Toiletten an der Uni) u.a. mit dem Alle-Frauen*-Referat
    zusammen, macht ein Alternativprogramm zur EM, veranstaltet ein
    OpenAir-Kino und wird sich noch stärker mit der Exzellenzinitiative
    beschäftigen. Auch Rassismus in der Lehre ist ein Thema.

    \textbf{HoPo/Öffentlichkeit/Kultur} (Geoffrey):

    Das "Open Raum"-Fest sei gelungen und eine gute Alternative zum früheren
    kommerziellen Campus Open Air. Die heutige Alternativ-Veranstaltung zur EM
    finde leider parallel zur StuPa statt. Man will die Stellwände auf dem Campus
    für studentische/politische Zwecke zurückerobern.

    \textbf{Soziales, Ökologie, Soziale Bewegung} (Artur):
    \begin{itemize}
        \item Am 18. Juni kommen Kollegen/Genossen aus Frankreich zur Uni, um
        von den dortigen Streiks und Auseinandersetzungen zu berichten. In Frankreich
        gibt es eine Studierendengewerkschaft, die sehr aktiv bei den Protesten
        ist.
        \item Am 23. Juni tagt der Wissenschaftsausschuss im Studierendenwerk,
        die Abgeordneten werden auch das Studierendenwerk begehen. Inhaltlich
        geht es u.a. um die Finanzierung und die Bedarfe des Studierendenwerks.
        Der AStA will u.a. der Anfrage der FDP mit dem Tenor, das
        Studierendenwerk sei bürokratisch und teuer, entgegenwirken.
    \end{itemize}

    \textbf{Ausschuss gegen Rechts} (Oliver und Ruben):

    Bei der Planung des "Tags gegen Rechts" sei der Rücklauf gering, daher wird
    er in der konzipierten Größenordnung nicht stattfinden. Es wird stattdessen
    am 5. Juli eine Veranstaltung gegen Rechts geben.

    \subsubsection{Anfragen an den AStA}

    Es gibt ergänzende Anmerkungen zur Begehung des Wissenschaftsrats: Der
    Darstellungsdruck sei sehr groß und viele sagten, alles sei toll, obwohl es
    nicht stimme. Dies zeige, dass Konkurrenz antidemokratisch sei. Es gehe dabei
    nicht um Erkenntnisse, sondern um Gläubigkeit.

    Der AStA will eine Einschätzung und Stellungnahme zur nächsten StuPa-Sitzung
    vorbereiten.

    \subsubsection{Dringlichkeitsanträge des AStA}

    Keine.

    \subsubsection{Aktuelle Stunde}

    Aus aktuellem Anlass wird die Exzellenzinitiative vorgeschlagen (sie wird
    zur Stunde in der Ministerpräsidentenkonferenz verhandelt).

    Franziska liest eine Meldung aus Berlin vor: Ein Kompromissvorschlag liege
    auf dem Tisch. Olaf Scholz sei für die ursprüngliche Fassung mit einigen
    Zugeständnissen gewesen, es solle jedoch weitere Änderungen geben.

    Es wird verschiedentlich eingeschätzt: Der jetzt vorgelegte Plan sei
    ambivalent, denn wenn alle Elite werden, gebe es keine Elite mehr. Eine
    Einigung wäre ärgerlich, aber die Kritik wirke und man solle auf jeden Fall
    weitermachen. Der Kampf lohne sich.

    Man müsse sich auch speziell gegen den Druck und die Erpressungsversuche
    wenden. Aus dem AS wird berichtet: Uni-Präsident Lenzen sagt auf Nachfrage,
    aus dem Hamsterrad auszutreten sei besser, weil es nicht so anstrengend sei
    und man Nützlicheres machen könne. Allerdings bekomme man dann kein Futter.

    Einige betonen die inhaltliche Seite der Kampagne: Was in Hamburg in den
    Clustern passiere, sei nicht nachhaltig. Es gehe darum, alle ökonomischen,
    ökologischen und sozialen Bedingungen für gesellschaftliche Nachhaltigkeit
    zu schaffen. Man solle nicht nur auf Exzellenz zuspitzen, sondern auf den
    Haushalt und die Grundfinanzierung der Hochschulen. Außerdem solle man auch
    auf die Professor*innen zugehen, das habe größere Wirkung.

    \textit{Zeitansage: 19.30 Uhr}

    \subsubsection{Feststellung der endgültigen Fassung des Teils B der Tagesordnung}

    Antrag Artur: Beitritt fzs vorziehen zu TOP 1: Mit m:1:3 angenommen.

    Vorschlag Präsidium: TOP 4 "Wahl-o-mat" vertagen, da Antragstellerin nicht da.
    Gegenantrag von Geoffrey auf Nichtbefassung mit m:3:4 angenommen.

    Vorschlag Präsidium: TOP 5 "Wahl im Sommersemester" vertagen, weil
    Antragssteller nicht da. Mit 11:8:7 angenommen.

    Die endgültige Fassung des Teils B der Tagesordnung wird einstimmig bei drei
    Enthaltungen festgestellt.

    %%% BESCHLOSSENE TAGESORDNUNG %%%
    \underline{Endgültige Fassung der Tagesordnung}
    \begin{enumerate}[label={\textbf{Top \theenumi}},leftmargin=*]
        % folgende Zeile bei Sitzungsfortsetzungen einkommentieren und x durch
        % die Zahl des zuletzt behandelten TOPs ersetzen
        %\setcounter{enumi}{x}
        \item Antrag auf Beitritt zum fzs (V1617-027)
        \item
            \begin{enumerate}
                \item Haushalt 2016/2017 1. Lesung (V1617-025) (90 Min.)
                \item Nachtragshaushalt 2015/2016 1. Lesung (V1617-026) (30 Min.)
            \end{enumerate}
        \item Akkreditierung (V1617-024) (60 Min.)
        \item Verschiedenes
    \end{enumerate}

    \subsubsection{Feststellung der Beschlussfähigkeit}

    Das Studierendenparlament ist mit 28 anwesenden Mitgliedern beschlussfähig.

    \subsubsection{Genehmigung der Protokolle der vorangegangenen Sitzungen}

    Es waren keine Protokolle zu genehmigen.

    \textit{Präsidiumspause von 10 Minuten bis 20.00 Uhr}

    \newpage
    %\vspace{0.5cm}
    {\Large \underline{Teil B}}

    % Teil B beginnt
    % Dieser wird sich für jede Sitzung unterscheiden.
    % Jeder TOP ist eine subsection.
    % Vorhandene Beispiele anfangs entfernen und durch tatsächliche TOPs ersetzen

    % Folgende Zeile bei Sitzungsfortsetzungen einkommentieren und x mit
    % der Zahl des zuletzt behandelten TOPs ersetzen.
    %\setcounter{subsection}{x}

    \textit{Jim übernimmt die Sitzungsleitung und Melf die Redeliste}

    \subsection{Antrag auf Beitritt zum fzs}

    Phillip stellt den Antrag vor: Der AStA solle von Fördermitglied zum
    richtigen Mitglied werden, das kostet der VS 30000 Euro im Jahr. Er geht
    kurz auf die Geschichte des fzs ein: Bis 1990 habe es den Vorläufer
    "Bund Deutscher Studierendenschaften" gegeben, der fzs wurde 1993 gegründet.
    Er hat 1,1 Mio Studierende als Mitglied. Der AStA Uni Hamburg wäre eine der
    größten Unis. Er finanziert eine Geschäftsstelle und vier Vorstände,
    Bahnfahrten, Seminare, Ausschüsse usw. Viele Unis können nicht so viel
    bezahlen. Es gibt ein gestaffeltes Stimmrecht mit 3 bis 6 Stimmen pro Uni.
    Eine Mitgliedschaft bringe neben den Kosten auch viel Arbeit mit sich. Ist
    es das wert? Mit einer Mitgliedschaft im fzs könne man mehr Dynamik z.B.
    gegen die Exzellenzinitiative entfachen und es gebe ein großes Potential an
    Mitstreiter*innen.

    Es gibt eine ausführliche, kontroverse Diskussion. Hauptsächlich aus
    Regenbogen/AL werden Vorbehalte und Fragen geäußert: Was passiert, wenn man
    austreten möchte? Wie kann man austreten? Was passiert mit dem vielen Geld?
    Kann man nicht eine Art Probezeit geben?

    Es wird erwidert: Es sei prinzipiell richtig, sich zu organisieren. Es bedeute
    nicht nur ein Mehraufwand, sondern auch Erleichterung der politischen Arbeit.
    Im fzs gebe es zur Zeit viel Internationalismus (Grüne) und Runde Tische
    (Sozialdemokraten), was fehlt, seien bundesweite Kampagnen. Das könne ein
    Beitrag aus Hamburg sein. Die Mitgliedshaft ende mit Austritt.

    \textit{Fraktionspause Alternative Linke bis 21.10 Uhr}

    Als Hauptsorge werden nochmal die 30.000 Euro im Jahr benannt.

    Es wird argumentiert, das sei für uns als größte Uni nicht so viel Geld im
    Verhältnis zu dem, was wir politisch bewegen können. Es sei auch nicht viel
    im Verhältnis zu den anderen Ausgaben im Haushalt der VS. Ein Eintritt wäre
    auch eine Selbstverpflichtung für den AStA, sich für sinnvolle Sachen
    einzusetzen, statt kommentierend am Rand zu stehen. Es gehe auch darum, ob
    der AStA sich als eine Ansammlung von Individuen begreift oder als ein
    politisches Gesamtgremium.

    Der Antrag, dem fsz beizutreten, wird mit m:1:5 angenommen.

    \subsection{Haushalt}

    \subsubsection{Haushalt 16/17 1. Lesung}

    \textit{Zeitansage: 21.50 Uhr}

    Timo und Philipp stellen den Haushaltsentwurf vor. Die Hauptveränderungen im
    Verhältnis zu vergangenen Jahren sind: Das kommerzielle "Campus Open Air"
    fällt weg, der Beitritt zum fzs ist beschlossen, die Budgets der
    Fachschaftsräte werden erhöht nach Studierendenzahl. Fast alle FSRe bekommen
    mehr, alle sind miteinander deckungsfähig. Globaltöpfe statt vieler kleiner
    Töpfe hätten sich bewährt.

    Auf die Frage nach dem Grundansatz dieses Haushalts antwortet Timo, dass es
    einen Streit gibt, ob es mehr Personal/Projektkräfte geben soll oder nicht.
    Die Ausgaben für Projektkräfte sind im Haushaltsentwurf erhöht, die Frage
    sei aber immer noch kontrovers.

    Melf hält als Vorsitzender des Haushaltsausschusses den Bericht vom
    Ausschuss: Viele Unklarheiten seien diskutiert und teilweise geklärt worden.

    In der Debatte wird darauf hingewiesen, dass die Budgets der FSRe früher
    nicht von den Studierendenanzahlen abhängig waren. Rechte ASten hätten es
    durchgesetzt und der Ansatz, da rauszukommen, sei richtig. Es wird angeregt,
    die (wenigen) FSRe, die durch die neue Formel weniger erhalten, zu fragen,
    ob ihre Arbeit dadurch eingeschränkt würde. Wenn ja, sollten sie aus der
    Formel rausgenommen werden.

    Es wird begrüßt, dass im Haushalt mehr Geld für Bundesfachschaftentagungen
    in Hamburg vorgesehen ist.

    Ferner wird erwähnt, dass es früher nicht nur die FSRK als teilautonomes
    Referat im AStA gab, sondern auch Fahrtkosten für verschiedene Aktivitäten,
    z.B. BuFaTas. Es wird vorgeschlagen, den Topf zu trennen in Fahrten zu und
    Ausrichtung von Tagungen.

    Es wird gefragt, was das politische Konzept hinter den 18 000 Euro mehr für
    Projektkräfte sei. Dies sei eine Verschiebung weg von politisch
    verantwortlichen AStA-Referent*innen und konträr zur politischen Entwicklung
    des AStA sowie zur neuen Referatsstruktur. Projektkräfte seien nur für einen
    Teilbereich zuständig und müssten nicht an AStA-Sitzungen teilnehmen. Es sei
    absurd, wenn AStA-Referent*innen und -Honorarkräfte nicht mal auf einer
    studentischen Vollversammlung oder zentralen Demo auftauchten.

    \textit{Zeitansage: 22.25 Uhr}

    Es entsteht eine Kontroverse zum Thema Projektkräfte.

    Auf der einen Seite wird vertreten, dies sei eine niedrigschwellige
    Möglichkeit, Leute für die politische Arbeit oder auch nur Projekte zu
    motivieren. Es wird sich dagegen verwehrt, dass die Projektkräfte weniger
    verantwortliche seien als Referent*innen.

    Auf der anderen Seite wird kritisiert, dass der AStA als ein Ort der Hilfe
    zur Selbstverwirklichung gesehen werde und gemeint, alle müssten sich mehr
    des Gesamts annehmen. Die Aussage dieses Haushalts sei: Wir wollen mehr
    Leute, die sich nicht so verantwortlich fühlen. Es wird dabei betont, dass
    es sich um ein strukturelles Problem handelt und niemand persönlich gemeint
    sei.

    \textit{GO-Antrag Melf: Verlängerung der Sitzung bis nach Verschiedenes.
    Kein Widerspruch}

    Als Mittlerposition wird vertreten: Viele sollten in der Referatsstruktur
    mitwirken können und nichts spreche dagegen, jemanden später als
    Referent*in zu bestätigen. Die Frage der Gesamtverantwortung sei ein
    Problem, aber keine Strukturfrage und unabhängig vom Status zu klären.

    \textit{Präsidiumspause von 10 Minuten bis 23.45 Uhr}

    Abstimmung: Mit 24:0:0 wird der Haushalt in die zweite Lesung überwiesen.
    Mit 24 abstimmenden Parlamentarier*innen ist das Parlament weiterhin
    eindeutig beschlussfähig.

    \subsubsection{Nachtragshaushalt 15/16 1. Lesung}

    Timo als Finanzreferent legt dar, dass der laufende Haushalt mit dem
    Nachtragshaushalt an die realen Ausgaben angepasst wurde.

    Melf als Haushaltsausschussvorsitzender erklärt, im Ausschuss seien sie den
    Nachtragshaushalt Punkt für Punkt durchgegangen.

    Nach einigen Fragen und Anregungen wird der Nachtragshaushalt mit 24:0:0 in
    die zweite Lesung überwiesen.

    \subsection{Akkreditierung}

    Phillip stellt den Antrag vor: Die Akkreditierung sei früher schlimmer gewesen,
    es sei einiges relativiert worden, aber das System sei dennoch grundsätzlich
    abzulehnen.

    In einer kurzen Aussprache wird vertreten, man müsse vor allem klären, wo
    man mit der Kritik, eigenen Aktivitäten und der Studienreform positiv
    hinwolle.

    In dem vorgelegten Text werden einmütig zwei Änderungen im zweiten Absatz
    vorgenommen: Das Wort "langwierigen" wird gestrichen und statt
    "Ministerialbürokratie" soll es "Ministerien" heißen.

    Abstimmung: Der Antrag wird mehrheitlich angenommen.

    \subsection{Verschiedenes}

    Das Präsidium teilt mit: Wer sich an der Vorbereitung der Diskussion mit
    Frau Fegebank am 30. Juni beteiligen möchte, kann am Donnerstag der Sitzung
    um 12 Uhr ins Präsidium kommen.

    \textit{Jim schließt um 0.50 Uhr die Sitzung.}

\end{document}
